\documentclass[../../../include/open-logic-section]{subfiles}

\begin{document}

\olfileid{his}{set}{pathology}
\olsection{病态}

然而，\emph{极限}的定义结果却容许了一些相当“病态的”构造。 

19世纪30年代左右，Bolzano（波尔查诺）发现了一个
\emph{处处连续}，但\emph{无处可导}的函数。
（遗憾的是，Bolzano 从未发表这一发现；数学家们直到1872年才首次接触到这个想法，这要归功于 Weierstrass（魏尔斯特拉斯）对同一想法的独立发现。）\footnote{这段历史记载在维基百科关于
\href{http://en.wikipedia.org/wiki/Weierstrass_function}{魏尔斯特拉斯函数}的文章的极其详尽的脚注中。}退一步说，这也相当令人惊讶。
很容易找到诸如 $|x|$ 这样的函数，它处处连续但在特定点不可导。
但是，一个处处连续却又\emph{无处}可导的函数则是完全不同的“怪物”。
试想一下，你会如何尝试画出这样一个函数。
为了确保其连续性，你必须能够一笔画成，笔尖永不离开纸面；
但为了确保其无处可导，你又必须不断突然改变运笔方向。

进一步的“病态”接踵而至。1874年1月5日，Cantor（康托尔）写信给 Dedekind（戴德金），提出了这样一个问题：

\begin{quote}
一个曲面（例如含边界的正方形）能否与一条直线（例如含端点的直线）一一对应，使得曲面上每一点都对应直线上一点，反之直线上每一点也对应曲面上一点？

我此刻仍然觉得，这个问题非常难以回答——尽管人们在此也极力想说\emph{不能}，以至于会认为证明几乎是多余的。[引自 \citealt{Gouvea2011}]
\end{quote}

但是，1877年，Cantor 证明他错了。事实上，一条直线和一个正方形所含的点数完全相同。他于1877年6月29日写信给 Dedekind 说：“\emph{je le vois, mais je ne le crois pas}”；即“我看见了，但我不相信”。在数学的“公认历史”中，这常常被用来表明这些新结果对当时的数学家来说\emph{从字面上就令人难以置信}。（相关通信见 \citet{Gouvea2011}，我们将在 \olref[his][set][mythology]{sec} 中再次提及。Cantor 证明的概览见 \olref[his][set][cantorplane]{sec}。）

受 Cantor 结果的启发，Peano（皮亚诺）开始思考是否可能将一条线\emph{光滑地}映射到一个平面上。这将是一条\emph{填满空间的曲线}。在 \citeyear{Peano1890}，Peano 恰恰构造出了这样一条曲线。这确实是反直觉的：Euclid（欧几里得）曾将线定义为“没有宽度的长度”（第 I 卷，定义 2），但 Peano 却表明，通过适当地卷曲一条线，其长度可以转化为宽度。
在 \citeyear{Hilbert1891}，Hilbert（希尔伯特）描述了一条稍具直观性的空间填充曲线，并附带了一些图示。该曲线是逐次构造的，以下是其构造的前六个阶段：

\begin{center}
\begin{tikzpicture}
\begin{scope}[xshift=20pt, yshift=20pt]
	\draw [oldiagcolorC, l-system={Hilbert curve, step=40pt, angle=90, axiom=L, order=1}] lindenmayer system;
\end{scope}
\begin{scope}[xshift=110pt, yshift=10pt]
	\draw [oldiagcolorC, l-system={Hilbert curve, step=20pt, angle=90, axiom=L, order=2}] lindenmayer system;
\end{scope}
\begin{scope}[xshift=205pt, yshift=5pt]
	\draw [oldiagcolorC, l-system={Hilbert curve, step=10pt, angle=90, axiom=L, order=3}] lindenmayer system;
\end{scope}
\begin{scope}[xshift=2.5pt, yshift=-97.5pt]
	\draw [oldiagcolorC, l-system={Hilbert curve, step=5pt, angle=90, axiom=L, order=4}] lindenmayer system;
\end{scope}
\begin{scope}[xshift=101.25pt, yshift=-98.75pt]
	\draw [oldiagcolorC, l-system={Hilbert curve, step=2.5pt, angle=90, axiom=L, order=5}] lindenmayer system;
\end{scope}
\begin{scope}[xshift=200.625pt, yshift=-99.375pt]
	\draw [oldiagcolorC, l-system={Hilbert curve, step=1.25pt, angle=90, axiom=L, order=6}] lindenmayer system;
\end{scope}
\draw[gray] (0pt,0pt) rectangle (80pt, 80pt);
\draw[gray] (100pt,0pt) rectangle (180pt, 80pt);
\draw[gray] (200pt,0pt) rectangle (280pt, 80pt);
\draw[gray] (0pt,-100pt) rectangle (80pt, -20pt);
\draw[gray] (100pt,-100pt) rectangle (180pt, -20pt);
\draw[gray] (200pt,-100pt) rectangle (280pt, -20pt);
\end{tikzpicture}  
\end{center}

在极限情形下——这个概念当时已有了严格定义——整个正方形会被实心红色填满。顺带一提，Hilbert 曲线处处连续却无处可导；直观上的原因是，在无穷极限下，函数每时每刻都在急剧改变方向。（我们将在 \olref[his][set][hilbertcurve]{sec} 更详细地概述 Hilbert 的构造。）

无论好坏，这些“病态的”几何构造都被当作怀疑几何直觉之可靠性的理由。它们几乎成了一种\emph{宣传}，用以推广一种新的做数学的方式，而这种方式最终在集合论中达到了顶峰。在该学科后来的神话塑造中，人们常说这些结果既完全严谨，又极其令人震惊。因此，它们起到了双重作用：既是对依赖几何直觉的警告，也是对新思维方式之丰硕成果的展示。

\end{document}